% =====================================================================
% Introduction — narrative body (precedes the FROZEN contributions list in
% introduction_contributions.tex). Scope-locked to frozen_claim_hierarchy_v1.md
% and OC-1..6. Order: problem -> evaluation entry -> interface limitation ->
% output-interface reframe -> three-proposition preview -> (contributions list).
% Wording strength <= §7.4/§7.6 and the abstract. No "comprehensive/全面".
% =====================================================================

\paragraph{Problem.}
Traffic disruptions---closures, congestion, lane reversals, staged recoveries---are announced and
updated in \emph{natural language}. A routing system that consumes such announcements must do two
things at once: \emph{ground} the language into the road graph (which segments are affected, in which
direction, for how long, and how severe), and then \emph{route} so that vehicles still reach their
destinations under the disruption. The useful objective is therefore not surface fidelity to the text
but the downstream traffic outcome: do trips \emph{complete}, and at what travel-time cost. We adopt
this as the evaluation entry point, reporting closed-loop completion first and travel time only
conditional on completion (\S\ref{sec:experiments}).

\paragraph{The interface gap in existing groundings.}
Most language-driven routing systems commit, implicitly, to one of two \emph{output interfaces} for the
grounding step. One emits a route directly from the description; the path is fluent but
\emph{uncheckable}---it carries no intermediate structure tying the decision to the specific disrupted
segments, and can be graph-invalid. The other emits a dense, per-edge cost or weight vector
($\Theta(|E|)$) for a downstream planner; this is verifiable but \emph{sequence-inefficient} and, as we
show, prone to collapsing into a fixed-capacity, scenario-invariant enumeration rather than a localized
edit. Both interfaces conflate \emph{what changed} with \emph{how to route}, and neither exposes a
compact, checkable representation of the disruption itself.

\paragraph{Our reframe: grounding as a choice of output interface.}
We treat constraint grounding as a first-class \emph{output-representation} decision and study it on a
fixed backbone, corpus, budget, and seed, varying \emph{only} how the grounded constraint is emitted.
We introduce \textbf{STAG-Route}, which grounds a natural-language constraint into \emph{typed sparse
anchors}: for each disrupted segment the model emits a small typed record
(\texttt{road$\mid$direction$\mid$range$\mid$type$\mid$level}), schema-validated and completed into the
graph by an external planner. The output is $\Theta(k)$ in the number of disrupted segments rather than
$\Theta(|E|)$, and every anchor is checkable against the road graph before it is acted on.

\paragraph{What we claim (and what we do not).}
We organize the evidence as three scoped propositions, each stated at the strength the experiments
support. \textbf{(P1) Representational adequacy.} On a constructed, method-agnostic, held-out
closed-loop SUMO stress test in which the unconstrained shortest path always violates the active
constraint, STAG-Route attains conditional constraint satisfaction $0.96$ at low excess detour and ranks
first of nine methods on both satisfaction and travel time, while a dense edge-weight grounder collapses
(CS $0.02$). The $0.96$ is robustness across the \emph{operational} stress buckets, not a
comprehensive-OOD claim: grounding degrades to $45$--$55\%$ under \emph{compositional} OOD (multiple
simultaneous conflicting events), a frontier we report and reconcile (\S\ref{sec:p4reconcile}).
\textbf{(P2) Interface robustness.} The typed-anchor \emph{interface} transfers across model families
(Qwen2.5, Llama-3.2)---schema-valid, no $\Theta(|E|)$ collapse---though absolute correctness is
family-dependent; we claim interface robustness, \emph{not} equal cross-family correctness, and make no
scale claim. \textbf{(P3) Topology transfer.} On held-out topologies the grounding localizes the
constraint to the correct \emph{road} (recall $0.93$) and the closed-loop advantage transfers to unseen
networks ($119/120$ vs.\ $0/120$ for the dense interface); edge-/segment-exact precision remains open. A separate mechanism finding---that even a near-oracle
numeric sensor does not close the gap---isolates constraint \emph{representation} from sensing
(\S\ref{sec:closedloop}). The contributions below make these precise.
